%%%%%%%%%%%%%%%%%%%%%%%%%%%%%%%%%%%%%%%%%%%%%%%%%%%%%%%%%%%%%%%%%%%%%%
% Lớp 10 | Đại số tổ hợp | Hoán vị. Chỉnh hợp. Tổ hợp
% Dạng 0D8?2-8: Hoán vị bàn tròn
% Số câu: 4  (sinh tự động bởi tach_cau.py; từ nay thầy sửa trực tiếp file này)
%%%%%%%%%%%%%%%%%%%%%%%%%%%%%%%%%%%%%%%%%%%%%%%%%%%%%%%%%%%%%%%%%%%%%%

%%%%%%%%%%%%%%%% ĐÚNG/SAI %%%%%%%%%%%%%%%%

% nguồn: MaHoaNH_DungSai_Lop10
\begin{ex}%[0D8T2-8]%[ma:549eb7ab]
	 Có $4$ nữ sinh tên là Huệ, Hồng, Lan, Hương và $4$ nam sinh tên là An, Bình, Hùng, Dũng cùng ngồi quanh một bàn tròn có $8$ chỗ ngồi. Chọn khẳng định đúng trong các khẳng định sau.
	 \choiceTF
		{Có $40320$ cách sắp xếp nam và nữ ngồi  tùy ý}
		{\True Có $5040$ cách sắp xếp nam và nữ ngồi tùy ý} 
		{Có $576$  cách sắp xếp biết nam và nữ ngồi xen kẽ nhau}
		{\True  Có $144$ cách sắp xếp nam và nữ ngồi xen kẽ nhau}
	\loigiai{ 
		\begin{itemize}
			\item Xếp nam và nữ ngồi tùy ý  có $7!=5040$ cách.
			\item Chọn $1$ bạn bất kì ngồi vào $1 $ vị trí ngẫu nhiên trên bàn tròn có $1$ cách, xếp $3$ bạn cùng giới tính còn lại vào  $3$ ghế có $3!$ cách, ếp $4$ bạn còn lại ngồi xen kẽ $4$ bạn đẫ xếp ở trên có $4!$ cách. Vậy có $3!\cdot 4!=144$ cách
		\end{itemize}
	} 
\end{ex}

% nguồn: MaHoaNH_DungSai_Lop10
\begin{ex}%[0D8T2-8]%[ma:26a0e199]
	Một hội nghị bàn tròn có các phái đoàn $3$ người  Anh, $5$ người Pháp và $7$ người Mỹ. Chọn khẳng định đúng trong các khẳng định sau.
	 \choiceTF
		 {\True Có $15!$ cách xếp các thành viên ngồi tùy ý}
		{Có $3!\cdot 4!\cdot 5!$ cách xếp các  thành viên ngồi tùy ý}
		{\True Có $7257600$ cách  xếp những người có cùng quốc tịch thì ngồi gần nhau}
		{Có $3628800$ cách xếp những người có cùng quốc  tịch thì ngồi gần nhau}
	\loigiai{ 
		\begin{itemize}
			\item Xếp các thành  viên ngồi tùy ý có $15!$ cách.
			\item Có $2!$ cách xếp $3$ phái đoàn vào bàn tròn. Với mỗi cách xếp thì có: $3!$ cách xếp các thành viên phái đoàn Anh, $5!$ cách xếp các thành viên phái đoàn Pháp, $7!$ cách xếp các thành viên phái đoàn Mỹ. Vậy có tất cả $2!\cdot 3!\cdot 5!\cdot 7! = 7257600$ cách xếp những người có cùng quốc tịch thì ngồi  gần nhau.
		\end{itemize}
	} 
\end{ex}


%%%%%%%%%%%%%%%% TỰ LUẬN %%%%%%%%%%%%%%%%

% nguồn: MaHoaNH_TuLuan_Lop10
\begin{bt}%[0D8H2-8]%[ma:fc63ea84]
	Có bao nhiêu cách sắp chỗ ngồi cho $6$ người vào $6$ ghế  xếp xung quanh một bàn tròn, nếu không có sự phân biệt giữa các ghế này?
	\loigiai{ 
		Vì bàn tròn không phân biệt đầu cuối nên để  xếp $6$ người ngồi quanh một bàn tròn ta cố định $1$ người và xếp $5$ người còn lại quanh người đã cố định. Vậy có $\mathrm{P}_5 = 5! = 120$ cách xếp.
	} 
\end{bt}

% nguồn: MaHoaNH_TuLuan_Lop10
\begin{bt}%[0D8H2-8]%[ma:34859e7a]
	\begin{enumerate}
		\item Hỏi có bao nhiêu cách xếp $6$ cặp vợ chồng ngồi xung quanh một chiếc bàn tròn, sao cho nam và nữ ngồi xen  kẽ nhau?
		\item Hỏi có bao nhiêu cách xếp $6$ cặp vợ chồng ngồi xung quanh một chiếc bàn  tròn, sao cho mỗi bà đều ngồi cạnh chồng của mình?
	\end{enumerate}
	 \loigiai{
		\begin{enumerate}
			\item Ta tiến hành xếp chỗ ngồi theo hai  công đoạn.\\
			Xếp $6$ nam ngồi quanh bàn tròn, có $(6-1)! = 5!$ cách xếp.\\ 
			Ta xem $6$ người nam vừa xếp là $6$ vách ngăn, vì $6$ người nam ngồi quanh bàn tròn nên có $6$ khoảng trống để xếp $6$ người  nữ nên có $6!$ cách xếp.\\
			Theo quy tắc nhân có $5! \cdot 6! = 86400$  cách.
			\item Ta tiến hành xếp chỗ ngồi theo hai  công đoạn.\\
			Xếp $6$ người chồng  ngồi quanh bàn tròn, có $(6-1)! = 5!$ cách xếp (vì vợ ngồi gần chồng).\\
			Mỗi cặp vợ chồng đổi chỗ cho nhau có  $2$ cách xếp mới nên có $2^6$ cách.\\
			Theo quy tắc nhân  có $5! \cdot 2^6 = 7680$ cách.
		\end{enumerate}
	 }
\end{bt}
